\section{Prolog Engine}
\subsection{Backchaining}
\definition{
Given a query $Q: \texttt{?- } A_1, \dots, A_n.$ and a rule $R: h \text{ :- } b_1, \dots, b_m.$, \defemph{backchaining} computes a new query by:
\begin{enumerate}
    \item finding a substitution that unifies the rule head $h$ with the selected goal $A_1$, and
    \item replacing $A_1$ in $Q$ with the body literals $b_1, \dots, b_m$, applying the substitution to all variables.
\end{enumerate}
}{}

\commandnote{
Backchaining motivates the terms \defemph{goals} and \defemph{subgoals}:
\begin{itemize}
    \item The literals in the query are ``goals'' that must be satisfied.
    \item Backchaining reduces a goal by replacing it with new ``subgoals'' from the matching rule body.
\end{itemize}
}

\definition{
The Prolog interpreter keeps backchaining from the \textbf{top} to the \textbf{bottom} of the program until the query
\begin{itemize}
    \item \defemph{succeeds}, i.e. contains no more goals, or 
    \item \defemph{fails}, i.e. backchaining becomes impossible
\end{itemize}
At this point, the interpreter issues an \defemph{answer}: \texttt{true} in the first case and \texttt{false} in the second.
}{}

\definition{
If a query contains variables, then Prolog will return an \defemph{answer substitution} as the answer to the query, i.e. the values for all the query variables accumulated during repeated backchaining.
}{}

\begin{example}[Backchaining]{ex:backchaining_trace}
Recall our knowledge base:
\begin{codeblock}
(1) human(socrates).
(2) human(leibniz).
(3) greek(socrates).
(4) fallible(X) :- human(X).
\end{codeblock}
Suppose we pose the query:
\begin{codeblock}
?- fallible(Y), greek(Y).
\end{codeblock}
The Prolog engine executes repeated \defemph{backchaining} steps:
\begin{enumerate}[label=\textbf{Step \arabic*:}, leftmargin=*]
    \item \textbf{Initial Query:} $Q_0 = \text{?- } \underline{\text{fallible}(Y)}, \text{greek}(Y).$
    \begin{itemize}
        \item The engine selects the first goal: $A_1 = \text{fallible}(Y)$.
        \item It scans the program from top to bottom and matches rule (4): $\text{fallible}(X) \text{ :- } \text{human}(X).$
        \item \textbf{Unification:} To make the head equal to $A_1$, it sets $[X \mapsto Y]$.
        \item \textbf{Replacement:} $A_1$ is replaced by the rule's body $\text{human}(Y)$.
    \end{itemize}
    This yields the new query:
    \[
    Q_1 = \text{?- } \underline{\text{human}(Y)}, \text{greek}(Y).
    \]
    \item \textbf{Resolving with a Fact:}
    \begin{itemize}
        \item The engine selects the new first goal: $A_1 = \text{human}(Y)$.
        \item Scanning from top to bottom, it finds fact (1): $\text{human}(\text{socrates}).$
        \item \textbf{Unification:} To make them equal, it binds $[Y \mapsto \text{socrates}]$.
        \item \textbf{Replacement:} Because a fact has an \textit{empty body}, $\text{human}(Y)$ is eliminated completely, and the binding is applied to all remaining goals.
    \end{itemize}
    This yields the simplified query:
    \[
    Q_2 = \text{?- } \underline{\text{greek}(\text{socrates})}.
    \]
    \item \textbf{Final Goal:}
    \begin{itemize}
        \item The engine selects $A_1 = \text{greek}(\text{socrates})$.
        \item It scans and finds fact (3): $\text{greek}(\text{socrates}).$
        \item Unification succeeds trivially (identical terms).
        \item The goal is eliminated.
    \end{itemize}
    This yields the empty query (denoted $\Box$ or \texttt{true}):
    \[
    Q_3 = \text{?- } \Box
    \]
\end{enumerate}
Since no subgoals remain, the query \textbf{succeeds}. Prolog outputs the accumulated variable binding as the \textbf{answer substitution}:
\begin{codeblock}
Y = socrates
\end{codeblock}
\end{example}

\begin{center}
\small
\begin{tabular}{cllll}
\toprule
\textbf{Step} & \textbf{Current Query} & \textbf{Selected Goal} & \textbf{Matched Clause} & \textbf{Substitution} \\
\midrule
$0$ & \texttt{?- fallible(Y), greek(Y).} & \texttt{fallible(Y)} & Rule (4) & $[X \mapsto Y]$ \\
$1$ & \texttt{?- human(Y), greek(Y).}    & \texttt{human(Y)}    & Fact (1) & $[Y \mapsto \text{socrates}]$ \\
$2$ & \texttt{?- greek(socrates).}       & \texttt{greek(socrates)} & Fact (3) & $\emptyset$ \\
$3$ & $\Box$ (\textit{empty query})      & ---                  & ---      & \textbf{Success: } $Y = \text{socrates}$ \\
\bottomrule
\end{tabular}
\end{center}

\subsection{Backtracking}
\definition[Prolog Search Procedure]{
The Prolog interpreter employs \defemph{top-down}, \defemph{left-to-right} \textit{Depth-First Search}. Concretely, the execution mechanism:
\begin{itemize}
    \item Resolves subgoals in \defemph{left-to-right} order.
    \item Matches the selected subgoal against the head literals of clauses in \defemph{top-to-bottom} order.
    \item If no clause matches, it fails and \defemph{backtracks} to the most recent \defemph{choice point} (or \defemph{backtrack point}).
    \item Otherwise, it backchains on the first matching clause and records remaining alternative matches as \defemph{choice points}.
\end{itemize}
We say that a goal $G$ matches a head $H$ iff they can be unified by binding variables to terms.
}{}

\commandnote{
We can force backtracking to compute more answers by typing \inlinecode{;}
}

\begin{example}[More Answers]{}
\begin{codeblock}
human(leibniz).
human(socrates).
human(plato).
greek(socrates).
greek(plato).
fallible(X) :- human(X).
\end{codeblock}

If we query \texttt{?- fallible(X), greek(X).}, Prolog yields \texttt{X = socrates}. Typing \texttt{;} prompts Prolog to backtrack, producing the next answer, \texttt{X = plato}.
\end{example}

\begin{example}[Backtracking and Choice Points]{ex:backtracking}
Consider the following knowledge base:
\begin{codeblock}
has_wheels(mytricycle, 3).
has_wheels(myrollerblade, 3).
has_wheels(mybmw, 4).
has_motor(mybmw).
car(X) :- has_wheels(X, 3), has_motor(X). % Rule 1
car(X) :- has_wheels(X, 4), has_motor(X). % Rule 2
\end{codeblock}
We pose the query to find an entity that is a car:
\begin{codeblock}
?- car(Y).
\end{codeblock}
The execution proceeds through depth-first search with \defemph{backtracking}:
\begin{enumerate}[label=\textbf{Step \arabic*:}, leftmargin=*]
    \item The goal is \texttt{car(Y)}. The engine finds a match with \textbf{Rule 1}. Because an untried alternative (\textbf{Rule 2}) exists, Prolog creates \textbf{Choice Point 1} and substitutes $[X \mapsto Y]$.
    \begin{itemize}
        \item Remaining goals: \texttt{has\_wheels(Y, 3), has\_motor(Y)}.
    \end{itemize}
    \item The first subgoal is \texttt{has\_wheels(Y, 3)}. The engine matches the first fact \texttt{has\_wheels(mytricycle, 3)}. Because an untried alternative fact exists (\texttt{myrollerblade}), Prolog creates \textbf{Choice Point 2} and binds $[Y \mapsto \text{mytricycle}]$.
    \begin{itemize}
        \item Remaining goal: \texttt{has\_motor(mytricycle)}.
    \end{itemize}
    \item The engine attempts to prove \texttt{has\_motor(mytricycle)}. Scanning the program from top to bottom reveals no matching facts or rules. The goal \textbf{fails}.
    \begin{itemize}
        \item The engine \textbf{backtracks} to the most recent choice point (\textbf{Choice Point 2}).
        \item The binding of $Y$ is \textbf{undone} ($Y$ becomes unbound again).
    \end{itemize}
    \item Prolog tries the next match for \texttt{has\_wheels(Y, 3)}, matching \texttt{has\_wheels(myrollerblade, 3)} with $[Y \mapsto \text{myrollerblade}]$. Since no further 3-wheel facts exist, Choice Point 2 is now exhausted.
    \begin{itemize}
        \item Remaining goal: \texttt{has\_motor(myrollerblade)}.
    \end{itemize}
    \item \texttt{has\_motor(myrollerblade)} finds no match and fails.
    \begin{itemize}
        \item Since Choice Point 2 is exhausted, Prolog backtracks further up to \textbf{Choice Point 1}.
        \item All bindings associated with Rule 1 are undone ($Y$ becomes unbound again).
    \end{itemize}
    \item Prolog now tries \textbf{Rule 2}: \texttt{car(X) :- has\_wheels(X, 4), has\_motor(X)}.
    \begin{itemize}
        \item Next subgoal: \texttt{has\_wheels(Y, 4)}.
        \item Matches \texttt{has\_wheels(mybmw, 4)} with binding $[Y \mapsto \text{mybmw}]$.
        \item Next subgoal: \texttt{has\_motor(mybmw)}. Matches the fact \texttt{has\_motor(mybmw)}.
    \end{itemize}
\end{enumerate}
No further goals remain ($\Box$). The query \textbf{succeeds}, and Prolog returns the answer substitution:
\begin{codeblock}
Y = mybmw
\end{codeblock}
\end{example}

\subsection{Functions}
\begin{excursion}[Remember - FOL Signature]{}
\definition{
A First-Order Logic Signature is a tuple $\Sigma := \tuple{\Sigma^f, \Sigma^p}$ where:
\begin{itemize}
\item $\Sigma^f := \bigcup_{k \in \mathbb{N}} \Sigma_k^f$ of \textbf{function constants} (or \textbf{terms}), where members of $\Sigma_k^f$ denote the $k$-ary functions on \textbf{individuals}.
\item $\Sigma^p := \bigcup_{k \in \mathbb{N}} \Sigma_k^p$ of \textbf{predicate constants}, where members of $\Sigma_k^p$ denote $k$-ary relations among individuals.
\item $\Sigma_k^f$ and $\Sigma_k^p$ are pairwise disjoint, countable sets of symbols for each $k \in \mathbb{N}$.
\item A $0$-ary \textbf{function constant} refers to a single \textbf{individual}; therefore, we call it an \textbf{individual constant}.
\end{itemize}
}{fol_sig}
\end{excursion}


In the semantics of First-Order Logic:
\begin{itemize}
    \item A \textbf{function constant} $f \in \Sigma_k^f$ is interpreted as a mapping $f^\mathcal{I}: D^k \to D$. A term like $\text{father}(\text{john})$ semantically translates to an \textbf{individual} in the domain (e.g., $\text{bob}$).
    \item A \textbf{predicate constant} $P \in \Sigma_k^p$ is interpreted as a relation $P^\mathcal{I} \subseteq D^k$. An atomic formula like $\text{loves}(\text{john}, \text{mary})$ evaluates strictly to a \textbf{truth value} depending on whether the tuple belongs to the relation.
\end{itemize}

In Prolog, however, the execution engine only evaluates \textbf{predicates} to determine truth or falsity. There is no mechanism where an expression evaluates to an individual, a term like \texttt{father(john)} will never semantically resolve to the individual \texttt{bob}.

To represent a function in Prolog, we must therefore encode it as a \textbf{predicate} by adding an extra argument for the output:
\begin{codeblock}
father(john, bob).    
\end{codeblock}
Here, \texttt{father(john, bob)} evaluates to \texttt{true}, while \texttt{father(john, alice)} evaluates to \texttt{false}. To find who the father of John is, we query:
\begin{codeblock}
?- father(john, X).
\end{codeblock}
where Prolog searches for the individual \texttt{X} that makes the predicate true.